\section{EBNF dan Konversi ke BNF}
\label{sec:ch08-ebnf}

\subsection{Notasi EBNF}

EBNF (\eng{Extended Backus-Naur Form}) menambahkan operator untuk
pengulangan, pilihan, dan pengelompokan:
\begin{table}[H]
\centering
\caption{Operator EBNF dan artinya}
\label{tab:ch08-ebnf}
\begin{tabularx}{\textwidth}{@{}l l X@{}}
\toprule
\textbf{Notasi} & \textbf{Arti} & \textbf{Contoh} \\
\midrule
\(X \mid Y\) & pilihan & \code{int | float} \\
\([X]\) & opsional & \code{[else stmt]} \\
\(\{X\}\) & nol atau lebih & \code{\{stmt\}} \\
\((X)\) & pengelompokan & \code{(a | b) c} \\
\bottomrule
\end{tabularx}
\end{table}

EBNF lebih ringkas dan mudah dibaca, tetapi parser generator seperti
Yacc/Bison memakai BNF. Karena itu, EBNF biasanya dikonversi ke BNF
sebelum implementasi.

\subsection{Aturan Konversi}

Konversi EBNF ke BNF dapat dilakukan secara sistematis:
\begin{enumerate}[label=\arabic*.,leftmargin=*]
  \item pilihan \(X \mid Y\) menjadi dua aturan terpisah;
  \item opsional \([X]\) menjadi \(X\) atau \(\varepsilon\);
  \item pengulangan \(\{X\}\) menjadi nonterminal bantu rekursif;
  \item pengelompokan \((X)\) diperlakukan sebagai nonterminal baru
        bila diperlukan.
\end{enumerate}

Sebagai contoh, aturan EBNF
\[
\text{params} = \text{param} \{\text{"," param}\}
\]
dikonversi menjadi
\[
\begin{aligned}
\text{params} &\to \text{param}\ \text{params\_tail},\\
\text{params\_tail} &\to \text{"," param}\ \text{params\_tail}
\mid \varepsilon.
\end{aligned}
\]
Aturan opsional
\[
\text{returnstmt} = \text{"return"}\ [\text{expr}]\ \text{";"}
\]
menjadi
\[
\text{returnstmt} \to \text{"return"}\ \text{opt\_expr}\ \text{";"},
\quad
\text{opt\_expr} \to \text{expr} \mid \varepsilon.
\]

\begin{figure}[H]
\centering
\begin{tikzpicture}[node distance=6mm and 10mm]
  \node[phase, text width=4.2cm] (ebnf) {EBNF\\{\{stmt\}}};
  \node[phase, text width=4.2cm, right=of ebnf] (bnf)
        {BNF\\stmt\_list};
  \node[phase, text width=4.2cm, below=of bnf] (rules)
        {stmt\_list \(\to\) stmt stmt\_list\\\(\mid \varepsilon\)};
  \draw[arrow] (ebnf) -- (bnf);
  \draw[arrow] (bnf) -- (rules);
\end{tikzpicture}
\caption{Konversi pengulangan EBNF menjadi nonterminal rekursif BNF}
\label{fig:ch08-ebnf-bnf}
\end{figure}

\subsection{Tata Bahasa Subset-C dalam EBNF}

Tata bahasa Subset-C ditulis dalam EBNF pada
\file{spec/grammar.ebnf}. Potongan awalnya ditunjukkan pada
\lstref{lst:ch08-ebnf}.

\includecode[
  language=EBNF,
  firstline=1,
  lastline=22,
  firstnumber=1,
  caption={Bagian awal \file{grammar.ebnf}},
  label={lst:ch08-ebnf}
]{../../code/subsetc/spec/grammar.ebnf}

Perhatikan penggunaan \(\{ \}\) untuk daftar pernyataan dan
\([ ]\) untuk bagian opsional. Notasi ini memudahkan pembacaan,
tetapi harus dikonversi sebelum dipakai oleh parser generator.

\subsection{Hasil Konversi ke BNF}

Hasil konversi lengkap ada pada \file{spec/grammar.bnf}. Potongannya
ditunjukkan pada \lstref{lst:ch08-bnf}.

\includecode[
  language=EBNF,
  firstline=1,
  lastline=24,
  firstnumber=1,
  caption={Bagian awal \file{grammar.bnf}},
  label={lst:ch08-bnf}
]{../../code/subsetc/spec/grammar.bnf}

Nama nonterminal bantu diberi akhiran \code{\_tail} atau awalan
\code{opt\_}. Konvensi ini memudahkan penelusuran asal aturan dan
membantu saat mentransformasi tata bahasa pada Bab IX.
